\section{怎样掌握文言的特殊词序}

古今汉语的词在句中的次序比较固定，
一般是主语在谓语之前，
动词谓语在宾语之前，
修饰语在被修饰语之前，
而且从古至今变化不大。
问题在于，
古代汉语有少数特殊的词序是现代汉语所没有的，
如果不了解这些特殊词序的用法，
就会给阅读古文带来一定的困难。

我们该怎样掌握这些特殊词序呢？

\textbf{1. 掌握宾语提前的语法规律。}

古代汉语中宾语提前的语法规律大致有三条，
我们应牢牢记住，
并注意在阅读古文时灵活运用。

第一条是：
疑问代词不论作动词的宾语，
还是作介词的宾语，
必须放在动词或介词的前面
（介词“以”的宾语，
即使不是疑问代词，
也常常放在“以”的前面）。

例如：

\noindent
\begin{enumerate}[label=\textcircled{\arabic*}]
	\item 沛公安在？ （沛公在哪里？） （《鸿门宴》）
	\item 卿欲何言？ （您想要说什么？） （《赤壁之战》）
	\item 缚者曷为者也？ （被绑着的人是干什么的呀？）（《晏子使楚》）
	\item 吾谁欺，欺天乎？ （我欺骗谁呢，欺骗老天吗？）（《论语·子罕》）
	\item 君之楚，将奚为北面？ （您到楚国去，又为什么向北方走呢？）（《战国策·南辕北辙》）
	\item 何以战？ （凭什么作战？）（《曹刿论战》）
	\item 微斯人，吾谁与归？ （不是这样的人，我跟谁一道走呢？）（《岳阳楼记》）
	\item 国胡以馈之？ （国家拿什么送给他们吃呢？）（《论积贮疏》）
	\item 吾孰与徐公美？ （我跟徐公比哪一个美呢？）（《邹忌讽齐王纳谏》）
	\item 余是以记之。 （我因此记下这件事。） （《石钟山记》）
\end{enumerate}

①②③④例句中的“安”、“何”、“曷”、“谁”都是疑问代词，
分别作动词“在”、“言”、“为”、“欺”的宾语，
位置都在动词之前。
⑤⑥⑦⑧例句中的“奚”、“何”、\\ “谁”、“胡”都是疑问代词，
分别作介词“为”、“以”、\\ “与”、“以”的宾语，
位置都在介词之前。
例句⑨中的“孰与”也是疑问代词放到了介词前面，
只是它已成为文言的一种固定格式，
经常连用，
意思是“跟······比哪一个······”。
例句⑩中的“是”是指代词，
作介词“以”的宾语，
位置也常放在“以”之前。
我们在译句时，
要把提前的宾语一律放到动词或者介词的后面去才行。

第二条是：
代词如果用作否定句中的动词的宾语，
一般就要放在动词之前，
并且紧接在否定词之后。

例如：

\noindent
\begin{enumerate}[label=\textcircled{\arabic*}]
	\item 自古及今，未之尝闻。（从古到今，不曾听说过这样的事。） （《论积贮疏》）
	\item 古之人不余欺也。（古代的人没有欺骗我们啊。） （《石钟山记》）
	\item 三岁贯女，莫我肯顾（多年来供养你，你丝毫不肯照看我。） （《硕鼠》）
	\item 每自比于管仲、乐毅，时人莫之许也。（常把自己跟管仲、乐毅相比，当时的人没有谁承认它。） （《隆中对》）
	\item 宁信度，无自信也。（宁可相信尺码，不相信自己的脚。） （《郑人买履》）
	\item 不患人之不已知，患不知人也。（不忧虑别人不了解自己，忧虑不了解别人啊。） （《论语·学而》）
\end{enumerate}

这六个例句中的“之”、“余”、“我”、“之”、“自”、\\ “已”，都是代词，
分别作各句动词“闻”、“欺”、“顾”、\\ “许”、“信”、“知”的前置宾语，
并且分别紧接在否定副词“未”、“不”、“莫”、“莫”、“无”、“不”之后。
我们在译这一类句子时，
要把提前的宾语一律放到动词后面去才行。

值得注意的是，
例句⑥中的“愚不知人也”并不把宾语“人”提到动词“知”的前面，
这是因为“人”是名词，
而不是代词。
这就跟前面的“不已知”的“已”是代词有所不同了。

第三条是：
用“之”、“是”、“之为”等作为宾语前置的标志，
它们的位置在前置宾语之后，
动词之前，
这种类型的词序有强调前置宾语的作用。

例如：

\noindent
\begin{enumerate}[label=\textcircled{\arabic*}]
	\item 譬若以肉投侯虎，何功之有哉？（就象将肉投到饿虎面前，有什么好效果呢？）（《信陵君窃符救赵》）
	\item 宋何罪之有？（宋国有什么罪呢？）（《公输》）
	\item 唯余马首是瞻。（只看我的马头。） （《冯婉贞》）
	\item 当臣之临流持竿时，心无杂虑，唯鱼之求。（当我面对河流手拿钓竿的时候，心里没有别的想法，只想钓到鱼。）（《列子·汤问》）
	\item 惟弈秋之为听。（只听取弈秋（的话）。）
\end{enumerate}
例句③④⑤的“唯（惟）······是（之）······”
或者
“唯（惟）······之（之为）······”
这类特殊的词序格式，
不仅强调了前置宾语，
而且表示了它的单一性、
排他性。
其中的“之”、\\ “是”、
“之为”等都已虚化成结构助词，
起到标志着宾语提前的作用，
不必译义。
我们在译句时，
只把前置宾语放到动词之后翻译就行了。

\textbf{2. 掌握定语后置的语法规律。}

定语在前，
为定语所修饰或限制的中心词在后。
这种词序，
不论在古代汉语里还是在现代汉语里，
基本如此。
但是，
古代汉语中有时可以把表示人或事物属性的修饰性定语放在它的中心词之后
（表示人或事物领属关系的限制性定语则不能），
这就需要我们认真辨析了。
当然，
我们在译句时，
还得按照现代汉语的习惯词序，
将后置的定语放到中心词之前进行翻译。

古代汉语的定语后置，
有三种主要形式。
第一种是“中心词\ +\ 定语\ +\ 者”，
第二种是“中心词\ +\ 之（或者“而”）\ +\ 定语\ +\ 者”，
第三种是“中心语\ +\ 之\ +\ 定语”。

例如：

\noindent
\begin{enumerate}[label=\textcircled{\arabic*}]
	\item 人马烧溺死者甚众。（烧死、淹死的人马很多。）《赤壁之战》
	\item 各处乡民来攻逆夷者，尚源源不绝。（各地前来攻杀侵略者的群众，还是源源不断。）《三元里抗英》
	\item 求人可使报秦者，未得。（找一个能够派去答复秦国的人，没有找到。）《廉颇蔺相如列传》
	\item 石之铿然有声者，所在皆是也。（铿锵地发出声音的石头，到处都这样。）
	\item 于是集谢庄少年之精技击者。（在这种情况下，召集了谢庄精通武术的年轻人。）《冯婉贞》
	\item 其石之突怒偃蹇，负土而出，争为奇状者，殆不可数。（那些突起似怒、屈曲起伏，背负泥土而耸出，争先恐后作出奇形怪状的石头，几乎多得数不清。）《钻钩潭西小丘记》
	\item 蚓无爪牙之利，筋骨之强······。（蚯蚓没有锐利的爪牙，强壮的筋骨······）《劝学》
	\item 居庙堂之高则忧其民，处江湖之远则忧其君。（处在高高的朝廷上就替百姓耽忧，处在偏远的江湖间就替国家耽忧。）《岳阳楼记》
\end{enumerate}

例句①②③属于第一种形式；
例句④⑤⑥属于第二种形式；
例句⑦⑧属于第三种形式。
其中的“之”和“者”，都是结构助词。
“之”起舒缓停顿作用，
不必译义；
“者”起煞尾作用，
相当于“的”。

此外，
古代汉语经常把数词放在名词后面，
而且不用量词。
我们在译句时，
应按现代汉语的词序把数词放到名词前面，
并加上恰当的量词。

例如：

\noindent
\begin{enumerate}[label=\textcircled{\arabic*}]
	\item 通计一舟，为人五。（总计整个核舟，刻了五个人。）《核舟记》
	\item 吏二缚一人诣王。（两个小官吏绑着一个人走到楚王面前来。）《晏子使楚》
	\item 马之千里者，一食或尽粟一石。（千里马，一顿有时能吃完一石小米。）《马说》
\end{enumerate}

\textbf{3. 掌握介词结构的翻译规律。}

文言的介词结构（或介宾结构）既可放在谓语前面作状语，
也可以放在谓语后面作补语。
其中由“与、为”组成的介词结构只能作状语，
译为现代汉语仍然是状语；
由“乎”组成的介词结构只能作补语，
译为现代汉语一般是状语。
需要注意的是“以”和“于”组成的介词结构。

由“以”组成的介词结构可以作状语，
也可以作补语，
译为现代汉语一般作状语用。

例如：
\newpage

\noindent
\begin{enumerate}[label=\textcircled{\arabic*}]
	\item 用花者，取花初敷时，用实者，成实时采。皆不可限以时月。（用花的药物，在花初开花时采；用果实的药物，在果实成熟时采。都不要用某时某月加以限制。）《采药》）
	\item 如君实责我以在位久，未能助上大有为，以膏泽斯民，则某知罪矣。（如果您拿在职时间已久，没有能帮助皇上有一番大作为，来给百姓造福责备我的话，那么我是愿意承认过错的。）（《答司马谏议书》）
\end{enumerate}

由“于”组成的介词结构主要作补语，
有时也作状语
（如果介词“于”表示“对于”的意思，
介词结构一定作状语），
译为现代汉语一般作状语用。
如“躬耕于南阳。”（亲自在南阳种地。）

总之，
介词结构在翻译时的位置一般放在谓语动词之前，
作状语用，
少数情况可以根据句意和语言习惯放在谓语动词后面翻译。

\textbf{4. 掌握谓语前置的翻译规律。}

为了强调谓语，
谓语也可提到主语前面，
古今汉语都有这种用法。
译句时，
我们要一律将谓语放到主语后面去。
例如：

\noindent
\begin{enumerate}[label=\textcircled{\arabic*}]
	\item 甚矣，汝之不惠！（你的傻气达到极点了！）（《愚公移山》）
	\item 安在公子能急人之困也！（公子您能够急解别人困境的表现在哪里呢？）（《信陵君窃符救赵》）
\end{enumerate}
\newpage